\begin{solution}{4}
设在某一时刻球壳形容器的电量为 $Q$。以液滴和容器为体系，考虑从一滴液滴从带电液滴产生器 $G$ 出口自由下落到容器口的过程。根据能量守恒有
\begin{equation}
m g h + k \frac{Qq}{h-R} = \frac{1}{2} m v^{2} + 2 m g R + k \frac{Qq}{R}.
\label{eq:1.s4}
\end{equation}
式中，$v$ 为液滴在容器口的速率，$k$ 是静电力常量。由此得液滴的动能为
\begin{equation}
\frac{1}{2} m v^{2} = m g (h-2R) - k \frac{Qq(h-2R)}{(h-R)R}.
\label{eq:2.s4}
\end{equation}
从上式可以看出，随着容器电量 $Q$ 的增加，落下的液滴在容器口的速率 $v$ 不断变小；当液滴在容器口的速率为零时，不能进入容器，容器的电量停止增加，容器达到最高电势。设容器的最大电量为 $Q_{\max}$，则有
\begin{equation}
m g (h-2R) - k \frac{Q_{\max} q (h-2R)}{(h-R)R} = 0.
\label{eq:3.s4}
\end{equation}
由此得
\begin{equation}
Q_{\max} = \frac{m g (h-R)R}{k q}.
\label{eq:4.s4}
\end{equation}
容器的最高电势为
\begin{equation}
V_{\max} = k \frac{Q_{\max}}{R}
\label{eq:5.s4}
\end{equation}
由\eqref{eq:4.s4}和\eqref{eq:5.s4}式得
\begin{equation}
V_{\max} = \frac{m g (h-R)}{q}
\label{eq:6.s4}
\end{equation}
\end{solution}